\section{Conclusion}

\mn{} connects differentiable routing to physical token compression within a
single local-to-latent forward path. Soft mass transport learns where visual
information should flow, an original-grid radius supplies a spatial prior
without fixing the merge assignments, and an exact-budget gather makes the
latent encoder operate on the same short sequence during training and
inference. The retained mass then parameterizes proportional latent attention;
a rank-one query/key augmentation provides an exact FlashAttention-compatible
implementation of that key bias.

The completed ablations support the spatial-routing component. Radius-3
routing improves the matched 150-epoch ImageNet controls and wins every
same-seed comparison against global, flattened, and degree-matched routing in
a three-seed CIFAR-100 campaign. At the full-model level,
\mn{} reaches 81.368\% top-1 with a $2\times$ carrier reduction, while the
matched-budget sweep records 81.360\% for \mn{} and 81.932\% for ToMe.
Full-grid local processing, recovery, and routing overhead mean that the
carrier ratio alone does not establish deployment speed.

The remaining submission evidence is narrow: one common-protocol numeric
packet for matched DTEM, latency, memory, and two checkpoint interventions;
and passive routing visualizations from the final ImageNet model. Until those
cells are filled, the supported claim is a trainable physical bottleneck with
a validated spatial-routing prior, rather than state-of-the-art deployment
performance.
